\section{Limitations}
\label{sec:limitations}

We state the limitations plainly; for an audit paper they are part of the contribution.
\begin{itemize}
\item \textbf{Provenance limit.} The pooled leakage$\to$reliance correlation is not recomputable
(per-fold leakage for two of three seeds was pruned); it is reproduced only at seed0 (sign) and carried
as \claim{FROZEN\_NOT\_RECOMPUTABLE} at pooled. Only alignment$\to$reliance is fully recomputed.
\item \textbf{Heterogeneity / small $n$.} The alignment$\to$reliance association is significant on 2a
but not on 2015; the alignment-vs-leakage contrast is at seed0, $n{=}42$, with non-significant
within-group partial betas.
\item \textbf{$\mathrm{align}_k$ is not a validated estimator}, only closer to reliance than raw
leakage.
\item \textbf{The RQ2 negative correlation is not a finding} --- it is a non-robust confound of
over-erasure and the random-$k$ anchor (Figure~\ref{fig:erasure}B); the RQ2 result is the absence of a
certified benefit (0/40).
\item \textbf{The natural-verification refit is minimal ERM}, reproducing the frozen deployment metrics
($|\Delta|{=}0.009$); it is an instrument to measure the branch-local ladder, not a tuned model, and its
conclusion (\claim{C16}) is a refusal, not a repair.
\item \textbf{The first-moment repair is construction-matched and narrow} (\claim{C12}): $73\%$ of the
pooled recovery is a mechanical identity (a first-moment aligner inverts a first-moment offset), the
effect is carried by one dataset and fails leave-one-dataset-out, and E4 overshoots the clean target ---
we report it as \textsc{strong within a controlled first-moment scope} and lead with the absolute
$+0.033$~bAcc gain, not the ratio.
\item \textbf{Second-moment repair is unresolved} (\claim{C14}, \claim{C17}): at the source-selected
operating point covariance-shrinkage does not beat a random-direction control, even oracle-directed; a
direction-specific advantage appears only in an injection-dominant regime we exclude. We do not claim
second-moment shortcuts are unconditionally unrepairable.
\item \textbf{Injected, not natural, controls.} The repair results use \emph{injected} positive controls
to test the machinery; they do not speak to natural or learned reliance, and must not be conflated with
the natural refusal (\S\ref{sec:rq4}).
\item \textbf{Head-only probes are mechanism tests, not natural-label-noise claims.} The two head-only
learned-reliance probes (\S\ref{sec:recoverability}) use controlled source-label manipulations on frozen
latents; they do \emph{not} prove that learned subject shortcuts cannot occur, and must not be read as a
natural-label-noise or ``resists weaponization'' claim. They show only that two cheap controlled
mechanisms (prevalence skew, subject-keyed task-conflict) do not license the paused PC2 refit or a
natural weaponization claim.
\item \textbf{Two datasets.} The binding robustness axis is leave-one-dataset-out; with two preset-ready
datasets it is under-powered (and 4F fails it). A \emph{full-backbone} learned-reliance refit (PC2) and a
$\geq 3$-dataset extension are future work, not current results; the head-only probes above are the cheap
mechanism substitute, and current evidence does not authorize the GPU refit.
\item \textbf{Proxies, not CMI.} L1/L2 use a label-conditional linear-probe posterior-KL surrogate and
a linear-probe advantage, not an unbiased mutual information; we say ``extractable conditional domain
information,'' not ``CMI.''
\item \textbf{Erasability $\neq$ full removal} (LEACE leaves a nonlinear residual); the one
target-unlabeled positive we cite (a non-CMI adaptation control) is seed0-only support, not an FSR
result.
\item \textbf{Scope.} Results are on specific backbones and two datasets; generality beyond them is not
claimed.
\end{itemize}
